\scp{\tsc{Tanimbar-Bomberai}}{11-05}
The main evidence for \tsc{Tanimbar-Bomberai} as a distinct node
within STB comes from a lexically specific split in the reflexes of *j.
Six words have *j > \it{r}
and six others *j > \it{s, r,} {\0}.
The words with *j > **r in all identified reflexes are shown
in \trf{11-22} with discussion following.
There are no certain reflexes of these words in Bedoanas
and Erokwanas, and only two likely reflexes in Arguni.
We discuss the reflexes of *j for those languages in more detail
in \srf{11-05-02}.

\begin{table}\tcp{\tsc{Tanimbar-Bomberai} *j > **r}{11-22}
\begin{adjustbox}{width=\textwidth}
	\begin{threeparttable}\st{2pt}
	\begin{tabular}{llllllll}\lsptoprule
*gloss	&	sun, day	&	nose	&	how m.?	&	gall bladder	&	navel	&	caterpillar\\
PMP	&	*qalə\tbr{j}aw	&	*(ŋ)i\tbr{j}uŋ	&	*pi\tbr{j}a	&	*qapə\tbr{j}u	&	*pusə\tbr{j}	&	*qulə\tbr{j}\\	\midrule
Yamdena	&	{\hr}le\tbr{r}-e	&	{\hr}i\tbr{r}i-ŋw	&	{\hr}fi\tbr{r}	&	{\hr}fe\tbr{r}u-	&	{\hr}fusa\tbr{r}	&	{\hr}ul-e\\	\midrule
Fordata	&	{\hr}le\tbr{r}a	&	{\hr}ni\tbr{r}u-n	&	{\hr}ifi\tbr{r}a	&	{\hr}fe\tbr{l}ur-n	&	{\hr}fua\tbr{r}	&	{\hr}ula\tbr{r}\\
Kei	&	{\hr}le\tbr{r}	&	{\hr}ni\tbr{r}u-n	&	{\hr}n-fi\tbr{r}	&	{\hr}fi\tbr{r}u-n	&	{\hr}fuha\tbr{r}	&	\cg{(laŋar)}\\
Onin	&	{\hr}re\tbr{r}a	&	{\hr}i\tbr{r}i-r	&	{\hr}fi\tbr{r}as	&	{\hr}ma\tbr{r}ik\su{†}	&	{\hr}fusa	&	\\
Sekar	&	{\hr}re\tbr{r}a	&	\cg{(idum)}	&	{\hr}fi\tbr{r}as	&	{\hr}ma\tbr{r}ik\su{†}	&	\cg{(ogoban)}	&	\\
Uruang.	&	{\hr}ne\tbr{r}a	&	{\hr}ni\tbr{r}iŋgain	&	{\hr}fi\tbr{r}a	&	\cg{(geram)}	&	{\hr}fusa\tbr{l}	&	\\
Arguni	&	\cg{(rúas)}	&	\cg{(sigyó)}	&	{\hr}fí\tbr{s} (loan?)	&	{\hr}ma\tbr{r}i\su{†}	&	\cg{(gomber)}	&	\\	\midrule
Bedoanas	&	\cg{(sáni)}	&	\cg{(suʔu)}	&		&		&		&	\\
Erokw.	&	\cg{(úruas)}	&	\cg{(suʔu)}	&		&		&		&	\\
Irarutu	&	{\hr}re, rre	&	\cg{(wegur}	&	\cg{(ɲa)}	&	-fi\tbr{r}- ‘bile’	&	{\hr}fwi\tbr{r}	&	\\
Kuri	&	{\hr}ro\tbr{r}, rəˈrɔː\tbr{r}	&	\cg{(iwɛ, ya-we-gə̆)}	&	\cg{(ɛnɔ)}	&		&	{\hr}ʷuɛ\tbr{r}ɛ	&	\\
		\lspbottomrule
	\end{tabular}
		\begin{tablenotes}\tnsize
			\item[†]	{Onin/Sekar \it{marik}, and Arguni \it{mari}
								all mean ‘bitter’ and are reflexes
								of intermediate **ma-pəju.}
			%\item[‡]	{The final part of Irarutu \it{wegur} nose may be connected
								%with the first part of PMP *(ŋ)ijuŋm though we consider this unlikely (\srf{11-05-0-02}).}
		\end{tablenotes}
	\end{threeparttable}
\end{adjustbox}
\end{table}

Regarding the reflexes of *qaləjaw,
\citet[272]{Jackson2014} gives both Irarutu \it{re} ‘day’
and \it{rre} [rəre] ‘sun,
day’, the latter of which he interprets as historically reduplicated
\it{r{\tl}re} with loss of the final syllable; *jaw.
Kuri \it{ror} retains *j > \it{r}
and alongside this word has [rəˈrɔːr] ‘daylight’.

Regarding the reflexes of *(ŋ)ijuŋ,
the Arguni, Bedoanas, and Erokwanas words for ‘nose’ are probably
not derivable from *ijuŋ as discussed in \srf{11-05-02-01}.
Linking Irarutu \it{wegur} [weŋgurə] to the doublet *ujuŋ or *ŋijuŋ
is possible, though much more difficult, on the assumption that initial \it{weg} can be analysed historically
as an independent morpheme -- as perhaps attested by Kuri \it{ya-weg} [ya-wegə].
See \srf{11-05-02-02} for more discussion.
Sekar \it{idum} is probably borrowed from Malay \it{idung} [iduŋ],
not inherited from *(ŋ)ijuŋ. Sekar also has \it{ipir} ‘nose’.

For the words in \trf{11-22} we can posit *j > **r at the level of Proto-Tanimbar-Bomberai.
This is the reflex in all languages for *qaləjaw ‘sun,
day’ and *(ŋ)ijuŋ ‘nose’ (excepting the potential Arguni,
Bedoanas, and Erokwanas reflexes).
Other words are not as straightforward.
Firstly, *j > \it{r} in *pija requires positing that Arguni \it{fís}
is a borrowing (perhaps from Biak,
Kowiai, or Geser \it{fis}).
Secondly, *j > **r in *qapəju requires positing medial **r >
\it{l} in Fordata \it{felur-n} ‘gall bladder’,
which finds independent support from the dissimilation of *d/*z
> \it{l} /{\gap}V\it{r} (see \trf{11-27} on \prf{tab:11-27}).
Thirdly, *j > **r for *pusəj requires positing irregular **r
> \it{l} in Uruangnirin \it{fusal} ‘navel’.
Lastly, *j > **r for *quləj requires sporadic loss of the final consonant in Yamdena.

All \tsc{Tanimbar-Bomberai} languages also have *R > \it{r},
with examples shown in \trf{11-23}.
Thus there is a partial merger of *j/*R > \it{r} in \tsc{Tanimbar-Bomberai}.
Several languages further have a merger of *R/*j with *d/*z (\srf{11-01-01}).
However, the \tsc{Bomberai Tip} languages (Onin,
Sekar, Uruangnirin) keep them distinct with *j > \it{r},
but *d/*z > \it{n}.\footnote{
		Word finally Yamdena has *d > \it{r}, \it{t}.
		Examples are: PMP *tuhu\tbr{d} > Yamdena \it{tu\tbr{r}-e}
		‘knee’ (Fordata \it{tura}, Kei \it{tuur}),
		*qabatə\tbr{d} > \it{bata\tbr{r}} ‘sago grub’,
		and *sumaŋə\tbr{d} > \it{smaŋa\tbr{t}} ‘soul’.
		Arguni \it{mesia simiat} ‘soul’ may be from *sumaŋəd with final
		*d > \it{t}, though this requires irregular *ŋ > {\0}.
		The latter is also in a semantic domain where loan
		and loan influences are common.}
Likewise, Yamdena keeps the reflexes partially distinct with *j > \it{r},
but *d/*z > \it{r, d}.

\begin{table}\tcp{\tsc{Tanimbar-Bomberai} *R > \it{r}}{11-23}
\begin{adjustbox}{width=\textwidth}
	\begin{threeparttable}\st{2pt}
	\begin{tabular}{llllllll}\lsptoprule
*gloss	&	house	&	new	&	blood	&	salt	&	thorn	&	stand	&	water\\
PMP	&	*\tbr{R}umaq	&	*baqə\tbr{R}u	&	*da\tbr{R}aq	&	*qasi\tbr{R}a	&	*du\tbr{R}i\su{‡}	&	*di\tbr{R}i	&	*wahiyə\tbr{R}\\	\midrule
Yamdena	&	\cg{(das)}	&	{\hr}be{\tl}be\tbr{r}y	&	{\hr}da\tbr{r}a-ny	&	{\hr}si\tbr{r}-e	&	{\hr}du\tbr{r}i	&	{\hr}na-mdi\tbr{r}y	&	{\hr}wey\\
Fordata	&	\cg{(rahan)}	&	\cg{(ŋorvaʔan)}	&	{\hr}la\tbr{r}a-n	&	{\hr}si\tbr{r}a	&	{\hr}lu\tbr{r}i-n	&	{\hr}n-di\tbr{r}i	&	{\hr}wea\tbr{r}\\
Kei	&	\cg{(raha)}	&	{\hr}vi\tbr{r}un	&	{\hr}la\tbr{r}	&	{\hr}si\tbr{r}	&	{\hr}lu\tbr{r}i-n	&	{\hr}di\tbr{r}	&	{\hr}wea\tbr{r}\\
Onin	&	{\hr}\tbr{r}uma	&	{\hr}per{\tl}pe\tbr{r}i	&	{\hr}ra\tbr{r}a	&	{\hr}si\tbr{r}a	&	{\hr}ru\tbr{r}i-r	&	{\hr}mari\tbr{r}i	&	{\hr}we\tbr{r}\\
Sekar	&	{\hr}\tbr{r}uma	&	{\hr}bìri{\tl}bí\tbr{r}i	&	{\hr}ra\tbr{r}a	&	{\hr}si\tbr{r}a	&	{\hr}ru\tbr{r}i-r	&	{\hr}mri\tbr{r}i	&	{\hr}we\tbr{r}\\
Uruang.	&	\cg{(kapala)}	&	{\hr}per{\tl}pe\tbr{r}i	&	{\hr}na\tbr{r}a	&	{\hr}si\tbr{r}a	&	{\hr}nu\tbr{r}i-n	&	{\hr}a-mri\tbr{r}i	&	{\hr}we\tbr{r}\\
Arguni	&	{\hr}\tbr{r}ume	&	{\hr}ir{\tl}í\tbr{r}	&	{\hr}ra\tbr{r}e	&	{\hr}si\tbr{r}er	&	\cg{(tór)}	&	\cg{(gáser)}	&	{\hr}wiyê\tbr{r}\\
Bedoanas	&	\cg{(windi)}	&	{\hr}ur{\tl}u\tbr{r}u	&	\cg{(rita)}	&		&	\cg{(tóʔor)}	&	\cg{(usar)}	&	{\hr}wie\tbr{r}\\
Erokw.	&	\cg{(widi)}	&	{\hr}pu{\tl}pwa\tbr{r}i	&	\cg{(rita)}	&	{\hr}si\tbr{r}ar	&	\cg{(toʔar)}	&	\cg{(waser)}	&	{\hr}we\tbr{r}\\
Irarutu	&	\cg{(san)}	&	\cg{(bunat)}	&	\cg{(wams)}	&	{\hr}te\tbr{r}ir\su{†}	&	{\hr}ru\tbr{r}	&	{\hr}n-mari\tbr{r}	&	{\hr}we\tbr{r}\\
Kuri	&	\cg{(wɛˈnɔm)}	&	\cg{(aˈbɔb)}	&	\cg{(wams)}	&		&	{\hr}ru\tbr{r}	&	{\hr}itə-bri\tbr{r}	&	{\hr}wɛ\tbr{r}\\
		\lspbottomrule
	\end{tabular}
		\begin{tablenotes}\tnsize
			\item[†]	{Irarutu \it{terir} either has irregular *s > \it{t} or is not from *qasiRa.}
			\item[‡]	{Yamdena,
								Onin, and Sekar reflexes mean ‘thorn, bone’.
								Other reflexes are only known to mean ‘bone’,
								though in some cases this could be due to lack
								of data.}
		\end{tablenotes}
	\end{threeparttable}
\end{adjustbox}
\end{table}

\begin{table}\tcp{\tsc{Tanimbar-Bomberai} *j > \it{s}, {\0}, \it{r}}{11-24}
%\begin{adjustbox}{width=\textwidth}
	\begin{threeparttable}\st{4pt}
	\begin{tabular}{llllllll}\lsptoprule
*gloss	&	rice	&	vine	&	name	&	foam	&	ySib	&	spike\\
PMP	&	*pa\tbr{j}ay	&	*waRə\tbr{j}\su{†}	&	*ŋa\tbr{j}an	&	*bu\tbr{j}əq	&	*hua\tbr{j}i\su{‡}	&	*su\tbr{j}a\su{Δ}\\	\midrule
Yamdena	&	{\hr}fa\tbr{s}-e	&	{\hr}wara\tbr{s}	&	{\hr}ŋa\tbr{r}-e	&	{\hr}bu{\tl}bu\tbr{s-}e	&	{\hr}wai-ny	&	{\hr}suan\\
Fordata	&	\cg{(wanat)}	&	{\hr}wara\tbr{t}	&	{\hr}na\tbr{r}a-n	&		&	{\hr}wa\tbr{r}i-n	&	{\hr}u\tbr{r}a\\
Kei	&	\cg{(kokat)}	&	{\hr}wara\tbr{t}	&	\cg{(mema-n)}	&		&	{\hr}wa\tbr{r}i-n	&	{\hr}hu\tbr{x}\\
Onin	&	{\hr}pa\tbr{s}a	&	{\hr}wara\tbr{s}	&	{\hr}ga\tbr{r}a-r	&		&	{\hr}wa\tbr{r}i-r	&	\\
Sekar	&	{\hr}pa\tbr{s}a	&	{\hr}wara\tbr{s}	&	{\hr}ga\tbr{r}a-r	&		&	{\hr}wa\tbr{r}e-r	&	\\
Uruang.	&	{\hr}fa\tbr{s}a	&	{\hr}wara\tbr{s}, wiri\tbr{s}	&	{\hr}(ŋ)ga\tbr{r}an	&		&	{\hr}wa\tbr{r}i-n	&	\\
Arguni	&	{\hr}pwá\tbr{s}en	&	{\hr}wári\tbr{r}	&	{\hr}gá\tbr{d}in	&		&	{\hr}twá\tbr{r}i	&	\\
Bedoanas	&	{\hr}pá\tbr{s}a	&	{\hr}wári\tbr{r}	&	{\hr}ga\tbr{d}iŋ	&		&		&	\\
Erokw.	&	{\hr}pá\tbr{s}a	&	{\hr}wári\tbr{r}	&	{\hr}ga\tbr{d}in	&		&		&	\\
Irarutu	&	{\hr}fa\tbr{s} 	&	{\hr}wara	&	\cg{(no)}	&	{\hr}fu\tbr{s}	&	\cg{(nefut)}	&	\\
		\lspbottomrule
	\end{tabular}
		\begin{tablenotes}\tnsize
			\item[†]	{The words in the languages of the Bomberai Peninsula were responses to ‘rope’.
								Uruangnirin \it{waras} is ‘rope’ and \it{wiris} ‘string’.}
			\item[‡]	{*huaji is reconstructed meaning ‘same sex younger sibling’.
								Onin, Sekar, and Arguni words were given in response to ‘friend’ (Indonesian \it{teman}).
								Arguni \it{tuari} is from *ta-huaji.
								(compare Arguni \it{támana} ‘father’ from *ta-ama.) Uruangnirin
								\it{wari-n} is ‘same sex sibling, friend’.}
			\item[Δ]	{*suja is reconstructed with the meaning ‘pitfall or trail spikes
								made of sharpened bamboo’.
								Kei \it{hux} ‘pointed bamboo stake’ is from Kei Besar with regular **r > \it{x}.}
		\end{tablenotes}
	\end{threeparttable}
%\end{adjustbox}
\end{table}

For another half a dozen words we cannot posit *j > **r in \tsc{Tanimbar-Bomberai}.
These words are shown in \trf{11-24}.

All reflexes of *pajay ‘rice’ have *j > \it{s}.
While we could propose lexically specific *j > \it{s},
initial *p = \it{p} in the Bomberai languages is irregular.
This indicates that this word has spread by contact to some extent,
particularly given that the diet in this region was traditionally sago based.
Words for ‘rice’ in the Papuan (non-Austronesian) languages of
the Bird’s Head are also frequently loans from an Austronesian
language ultimately reflecting *pajay.

The two identified reflexes of PMP *bujəq ‘foam’ in Yamdena
and Irarutu also show *j > \it{s.} \citet{BlustTrussel2020} derive
Yamdena \it{bu{\tl}bus-e} from the doublet *busa,
in which case medial \it{s} would be a shared inheritance -- though
PMP *busa may be spurious.\footnote{
		The non-Oceanic forms given by \citet{BlustTrussel2020} in support
		of *busa may be Malay loans; e.g.
		\citet{BlustTrussel2020} give Buru \it{busa} in support of this
		form, but this is a Malay loan,
		as shown by *b = \it{b.} Buru has native *bujəq > \it{fue-n.}}
The sound correspondences for the final consonant of *waRəj
are diverse, and some of these consonants (particularly \it{t}) may be reflexes
of a nominal suffix \it{-t} that is productive in many languages
of the region, but perhaps frozen in this particular etymon for many languages.

Reflexes of *ŋajan ‘name’ have *j > \it{d} in \tsc{\tsc{East
Bomberai}} and *j > \it{r} in other languages.
We cannot posit shared *j > **d for this word,
as it does not merge with \tsc{Bomberai Tip} *d/*z,
which have *d/*z > \it{n}.
Instead, we need to posit *j > \it{d} for \tsc{East Bomberai},
and *j > \it{r} for other languages.
Reflexes of *huaji and *suja have *j > {\0} in Yamdena,
but *j > **r in other languages,
meaning that we cannot straightforwardly posit a shared sound
change for these items at the level of Proto-Tanimbar-Bomberai.

To summarise, the evidence for \tsc{Tanimbar-Bomberai} comes
from a split in *j whereby six words merge *j/*R > **r.
Reflexes attesting *j > **r which are not clear and
unambiguous inheritances from PMP are currently lacking for some
\tsc{East Bomberai} languages meaning that their position
within \tsc{Tanimbar-Bomberai} and STB must
be reassessed if additional reflexes are identified.
